% ═══════════════════════════════════════════════════════════════
% LH-02a: Lehrerhinweise -- La familia (Grundwortschatz)
% Abgeleitet aus LE-02a (Score 9.3/10)
% ═══════════════════════════════════════════════════════════════

\documentclass[11pt, a4paper]{article}

% ═══════════════════════════════════════════════════════════════
% SW-MODUS TOGGLE
% ═══════════════════════════════════════════════════════════════
\newif\ifswmodus
\swmodusfalse    % Standard: Farbe

\usepackage[utf8]{inputenc}
\usepackage[T1]{fontenc}
\usepackage[ngerman]{babel}
\usepackage{helvet}
\renewcommand{\familydefault}{\sfdefault}

\usepackage[margin=2cm]{geometry}
\usepackage{enumitem}
\usepackage{tabularx}
\usepackage{booktabs}
\usepackage{longtable}

\usepackage{xcolor}
\usepackage{amssymb}
\usepackage{tcolorbox}
\tcbuselibrary{skins}

\usepackage{fancyhdr}
\usepackage{lastpage}
\usepackage{needspace}

% Farben - Basis
\definecolor{spanisch-color}{HTML}{c41e3a}
\definecolor{grau-color}{HTML}{888888}
\definecolor{hellgrau-color}{HTML}{f5f5f5}
\definecolor{basis-color}{HTML}{2a7a4b}
\definecolor{standard-color}{HTML}{2c5aa0}
\definecolor{erweiterung-color}{HTML}{7b1fa2}
\definecolor{loesung-color}{HTML}{c62828}

% SW-Farben
\definecolor{sw-dark}{HTML}{000000}
\definecolor{sw-gray}{HTML}{666666}
\definecolor{sw-white}{HTML}{ffffff}

% Bedingte Farbzuweisung
\ifswmodus
    \colorlet{spanisch}{sw-dark}
    \colorlet{grau}{sw-gray}
    \colorlet{hellgrau}{sw-white}
    \colorlet{basis}{sw-dark}
    \colorlet{standard}{sw-dark}
    \colorlet{erweiterung}{sw-dark}
    \colorlet{loesung}{sw-dark}
\else
    \colorlet{spanisch}{spanisch-color}
    \colorlet{grau}{grau-color}
    \colorlet{hellgrau}{hellgrau-color}
    \colorlet{basis}{basis-color}
    \colorlet{standard}{standard-color}
    \colorlet{erweiterung}{erweiterung-color}
    \colorlet{loesung}{loesung-color}
\fi

\newcommand{\fachfarbe}{spanisch}

% Variablen
\newcommand{\thema}{La familia -- Grundwortschatz}
\newcommand{\sequenz}{02}
\newcommand{\block}{a}
\newcommand{\fach}{Spanisch}
\newcommand{\klasse}{7}
\newcommand{\dauer}{90 Min. (Doppelstunde)}
\newcommand{\lerngruppengroesse}{24}
\newcommand{\datum}{\today}

% Kopf-/Fußzeile
\pagestyle{fancy}
\fancyhf{}
\renewcommand{\headrulewidth}{0.4pt}
\renewcommand{\footrulewidth}{0.4pt}
\lhead{\fach{} -- Klasse \klasse}
\chead{\textbf{LH-\sequenz\block{} (Lehrerhinweise)}}
\rhead{\datum}
\lfoot{}
\rfoot{Seite \thepage\ von \pageref{LastPage}}

% Befehle
\newcommand{\B}{\textcolor{basis}{\textbf{[B]}}}
\newcommand{\Std}{\textcolor{standard}{\textbf{[S]}}}
\newcommand{\E}{\textcolor{erweiterung}{\textbf{[E]}}}
\newcommand{\loesung}[1]{\textcolor{loesung}{\textbf{#1}}}
\newcommand{\es}[1]{\textcolor{\fachfarbe}{\textbf{#1}}}

% Boxen (SW-kompatibel: schwebender Titel)
\newtcolorbox{kurzplan}{
    colback=hellgrau,
    colframe=\fachfarbe,
    colbacktitle=white,
    coltitle=\fachfarbe,
    boxrule=1pt,
    arc=0pt,
    fonttitle=\bfseries,
    attach boxed title to top left={yshift=-2mm, xshift=5mm},
    boxed title style={boxrule=0pt, colback=white},
    title=Kurzplan
}

\newtcolorbox{differenzierung}{
    colback=white,
    colframe=grau,
    colbacktitle=white,
    coltitle=grau,
    boxrule=0.5pt,
    arc=0pt,
    fonttitle=\bfseries,
    attach boxed title to top left={yshift=-2mm, xshift=5mm},
    boxed title style={boxrule=0pt, colback=white},
    title=Differenzierung
}

\newtcolorbox{fehlerbox}{
    colback=white,
    colframe=loesung,
    colbacktitle=white,
    coltitle=loesung,
    boxrule=0.5pt,
    arc=0pt,
    fonttitle=\bfseries,
    attach boxed title to top left={yshift=-2mm, xshift=5mm},
    boxed title style={boxrule=0pt, colback=white},
    title=Häufige Fehler
}

\newtcolorbox{materialbox}{
    colback=white,
    colframe=\fachfarbe,
    colbacktitle=white,
    coltitle=\fachfarbe,
    boxrule=0.5pt,
    arc=0pt,
    fonttitle=\bfseries,
    attach boxed title to top left={yshift=-2mm, xshift=5mm},
    boxed title style={boxrule=0pt, colback=white},
    title=Material-Checkliste
}

% ═══════════════════════════════════════════════════════════════
\begin{document}

% Kopfbereich
\begin{center}
    {\Large\bfseries\textcolor{\fachfarbe}{Lehrerhinweise: \thema}}\\[0.3em]
    {\normalsize LH-\sequenz\block{} | \dauer}
\end{center}

\vspace{10pt}

% ═══════════════════════════════════════════════════════════════
% 1. KURZPLAN
% ═══════════════════════════════════════════════════════════════

\section*{1. Stundenverlauf (Kurzplan)}

\textbf{1. Stunde (45 Min.)}

\begin{tabularx}{\textwidth}{|c|l|X|l|}
\hline
\textbf{Zeit} & \textbf{Phase} & \textbf{Inhalt} & \textbf{Material} \\
\hline
0--5 & Einstieg & Familienfoto zeigen, ``¿Quién es?'' & Bild \\
\hline
5--15 & Input & 10 Familienbegriffe mit Bildkarten einführen & Bildkarten \\
\hline
15--25 & Übung 1 & Chorisches Sprechen: Aussprache üben & -- \\
\hline
25--35 & Übung 2 & Memory-Spiel in Gruppen (Bild-Wort) & Memory-Karten \\
\hline
35--40 & Sicherung & Schnelle Abfrage: Bild $\rightarrow$ Wort & Bildkarten \\
\hline
40--45 & Überleitung & Merkkasten im AB zeigen, beginnen & AB-02a \\
\hline
\end{tabularx}

\vspace{1em}

\textbf{2. Stunde (45 Min.)}

\begin{tabularx}{\textwidth}{|c|l|X|l|}
\hline
\textbf{Zeit} & \textbf{Phase} & \textbf{Inhalt} & \textbf{Material} \\
\hline
0--5 & Wiederholung & Vokabelquiz: ``Wie heißt Mutter?'' & -- \\
\hline
5--15 & Anwendung & Stammbaum an Tafel gemeinsam beschriften & Tafel, AB \\
\hline
15--25 & Dialog & Redemittel ``¿Quién es? -- Es mi...'' üben & PA-Karten \\
\hline
25--35 & Festigung & AB Aufgabe 3+4 bearbeiten & AB-02a \\
\hline
35--40 & Sicherung & 2-3 Dialoge vorführen lassen & -- \\
\hline
40--45 & Abschluss & Zusammenfassung, HA aufgeben & -- \\
\hline
\end{tabularx}

% ═══════════════════════════════════════════════════════════════
% 2. DIFFERENZIERUNG
% ═══════════════════════════════════════════════════════════════

\section*{2. Differenzierung}

\begin{differenzierung}
\textbf{Legende:} \B{} Basis (BR) | \Std{} Standard (MR) | \E{} Erweiterung (GY)

\vspace{0.5em}

\textbf{WICHTIG:} Diese Marker sind NUR für die Lehrkraft. Auf Schülermaterial erscheinen KEINE Niveaumarkierungen!
\end{differenzierung}

\vspace{1em}

\begin{tabularx}{\textwidth}{|l|X|X|X|}
\hline
& \B{} \textbf{Basis} & \Std{} \textbf{Standard} & \E{} \textbf{Erweiterung} \\
\hline
\textbf{Aufgabe 1} & Zuordnung mit Bildhilfe & Zuordnung selbstständig & + Sätze bilden \\
\hline
\textbf{Aufgabe 2} & Lückentext mit Wortliste & Lückentext (Wortliste optional) & + Eigene Sätze schreiben \\
\hline
\textbf{Aufgabe 3} & Stammbaum mit Vorgaben & Stammbaum selbstständig & Eigene Familie einzeichnen \\
\hline
\textbf{Aufgabe 4} & Dialog mit Vorlage & Dialog mit Stichworten & Freier Dialog (ohne Hilfe) \\
\hline
\end{tabularx}

\vspace{1em}

\textbf{Konkrete Hinweise:}
\begin{itemize}
    \item \B{} SuS mit LRS: Zeitzugabe (+10 Min.), vergrößerte AB-Version
    \item \B{} DaZ-SuS: Familienbegriffe deutsch-spanisch vorab besprechen, Bildkarten nutzen
    \item \E{} Leistungsstarke: Eigene Familie auf Spanisch beschreiben (3-4 Sätze)
    \item \E{} Schnelle SuS: Erweiterungswortschatz (los padres, los abuelos, el hijo/la hija)
\end{itemize}

% ═══════════════════════════════════════════════════════════════
% 3. HÄUFIGE FEHLER
% ═══════════════════════════════════════════════════════════════

\section*{3. Häufige Fehler}

\begin{fehlerbox}
\begin{tabularx}{\textwidth}{|l|X|X|}
\hline
\textbf{Fehler} & \textbf{Beispiel} & \textbf{Reaktion} \\
\hline
Aussprache \textit{hermano} & *[hermano] mit H & H ist stumm! Vorsprechen: [er'mano] \\
\hline
Verwechslung \textit{primo/prima} & ``primo'' für Cousine & Endungsregel: -o = m, -a = f \\
\hline
Artikel vergessen & *``padre se llama...'' & Immer mit Artikel lernen: \textit{el} padre \\
\hline
Falscher Artikel & *``la padre'' & Merkhilfe: Endung zeigt Genus! \\
\hline
Verwechslung \textit{tío/abuelo} & Onkel und Großvater & Stammbaum-Visualisierung nutzen \\
\hline
\end{tabularx}
\end{fehlerbox}

% ═══════════════════════════════════════════════════════════════
% 4. MATERIAL-CHECKLISTE
% ═══════════════════════════════════════════════════════════════

\section*{4. Material-Checkliste}

\begin{materialbox}
\begin{itemize}[label=$\square$]
    \item AB-02a.pdf (\lerngruppengroesse{} Kopien)
    \item ML-02a.pdf (1$\times$ für Lehrkraft)
    \item Lehrwerk: Qué pasa 1, S. 54-57
    \item Bildkarten Familienmitglieder (10 Karten, laminiert)
    \item Memory-Karten (6 Sets für Gruppenarbeit)
    \item Tafelmarker / Kreide (für Stammbaum)
    \item Optional: LP-02a.html (Tablets/PCs)
    \item Optional: Familienfoto (authentisch, aus Lehrwerk)
\end{itemize}
\end{materialbox}

% ═══════════════════════════════════════════════════════════════
% 5. LÖSUNGEN
% ═══════════════════════════════════════════════════════════════

\section*{5. Lösungen AB-02a}

\textbf{Merkkasten: Familienmitglieder}

\begin{tabular}{ll|ll}
\es{el padre} & \loesung{der Vater} & \es{la madre} & \loesung{die Mutter} \\
\es{el hermano} & \loesung{der Bruder} & \es{la hermana} & \loesung{die Schwester} \\
\es{el abuelo} & \loesung{der Großvater} & \es{la abuela} & \loesung{die Großmutter} \\
\es{el tío} & \loesung{der Onkel} & \es{la tía} & \loesung{die Tante} \\
\es{el primo} & \loesung{der Cousin} & \es{la prima} & \loesung{die Cousine} \\
\end{tabular}

\vspace{1em}

\textbf{Aufgabe 1: Zuordnung} (4 Punkte)

\begin{tabular}{ll}
\es{el padre} & $\rightarrow$ \loesung{der Vater} \\
\es{la madre} & $\rightarrow$ \loesung{die Mutter} \\
\es{el hermano} & $\rightarrow$ \loesung{der Bruder} \\
\es{la hermana} & $\rightarrow$ \loesung{die Schwester} \\
\end{tabular}

\vspace{1em}

\textbf{Aufgabe 2: Lückentext} (4 Punkte)

\begin{tabular}{ll}
a) & Mi \loesung{padre} se llama Carlos. \\
b) & Mi \loesung{madre} se llama María. \\
c) & Tengo un \loesung{hermano} y una \loesung{hermana}. \\
d) & Mi \loesung{abuelo} es muy simpático. \\
\end{tabular}

\vspace{1em}

\textbf{Aufgabe 3: Stammbaum} (6 Punkte)

\begin{tabular}{ll|ll}
1 & \loesung{el abuelo} & 2 & \loesung{la abuela} \\
3 & \loesung{el padre} & 4 & \loesung{la madre} \\
5 & \loesung{el tío} & 6 & \loesung{la tía} \\
7 & YO & 8 & \loesung{el hermano / la hermana} \\
9 & \loesung{el primo} & 10 & \loesung{la prima} \\
\end{tabular}

\vspace{1em}

\textbf{Aufgabe 4: Mini-Dialog} (6 Punkte)

Mögliche Antworten:
\begin{itemize}
    \item B1: \loesung{Es mi padre. Se llama [Name].}
    \item B2: \loesung{Es mi madre. Se llama [Name].}
    \item A3: \loesung{¿Tienes hermanos? / ¿Quién es él?}
    \item B3: \loesung{Sí, tengo un hermano. / Es mi hermano.}
\end{itemize}

\textit{Bewertung: Struktur ``Es mi...'' (2P), Familienbegriff (2P), Name (1P), Aussprache (1P)}

% ═══════════════════════════════════════════════════════════════
% 6. HAUSAUFGABE
% ═══════════════════════════════════════════════════════════════

\section*{6. Hausaufgabe}

\begin{itemize}
    \item Lehrwerk S. 56, Übung 2 (Stammbaum beschriften)
    \item Vokabeln lernen: 10 Familienmitglieder mit Artikel
    \item Optional: Lernpfad LP-02a.html (digital)
    \item Optional: Eigenen Familienstammbaum zeichnen (Vorbereitung 02b)
\end{itemize}

% ═══════════════════════════════════════════════════════════════
% 7. REFLEXIONSNOTIZEN
% ═══════════════════════════════════════════════════════════════

\section*{7. Reflexionsnotizen (nach Durchführung)}

\begin{tabularx}{\textwidth}{|l|X|}
\hline
\textbf{Was lief gut?} & \\[2em]
\hline
\textbf{Was lief nicht?} & \\[2em]
\hline
\textbf{Zeitmanagement} & \\[2em]
\hline
\textbf{Für nächstes Mal} & \\[2em]
\hline
\end{tabularx}

\end{document}
